\chapter{Кубиты, комплексные векторы состояния и наблюдаемые}

\lectureoverview{Спин одного электрона — первый подлинно квантовый бит. Его можно приготовить относительно любого пространственного направления, однако фиксированный детектор дает лишь два исхода. Мы увидим, где перестает работать классическая магнитная картина, а затем построим комплексное векторное пространство, правило вероятностей и матричный язык, необходимые для описания того, что в действительности происходит в эксперименте.}

\section{Один спин и два смысла слова «вектор»}

Классический бит имеет два возможных состояния, и ничего между ними. Мы можем назвать их \(0\) и \(1\), орлом и решкой или вверх и вниз; эти названия — метки двух различных точек в классическом пространстве состояний. Квантовый бит тоже дает две альтернативы при измерении фиксированным способом, но его нельзя описать всего двумя точками пространства состояний. Простейший пример — спин электрона.

В этом примере есть два разных понятия вектора, и если их путать, все дальнейшее становится излишне загадочным. Во-первых, есть обычный \emph{пространственный вектор}: стрелка в трехмерном пространстве с величиной и направлением. У электрона есть магнитный момент, который для нашей нынешней цели ведет себя как крошечный магнит. Его величина фиксирована, тогда как в процедуре приготовления ему можно приписать пространственное направление.

Второй объект — абстрактный вектор в комплексном векторном пространстве. Этот \emph{вектор состояния} будет описывать квантовое состояние. Это не просто магнитная стрелка, записанная в другой нотации. Связь между пространственными направлениями и векторами состояния реальна и важна, но это не тождество двух объектов. К этой связи мы придем постепенно.

\section{Где классическая картина перестает работать}

Начнем намеренно с классической истории. Поместим крошечный магнит во внешнее сильное магнитное поле. Если только он уже не ориентирован вдоль поля, он прецессирует вокруг него. По мере излучения энергии прецессия затухает, и магнит устанавливается в ориентации с наименьшей энергией. Выбирая направление приложенного поля, дожидаясь этой релаксации, а затем убирая поле, мы, по-видимому, можем приготовить магнитный момент в любом желаемом направлении.

Та же установка подсказывает способ детектирования. Включим поле в выбранном направлении и будем наблюдать излучение, возникающее, пока маленький магнит устанавливается. Классический магнит, выровненный вдоль поля, не должен излучать энергию. Магнит, первоначально направленный против поля, должен излучить максимальное количество энергии, а при промежуточном угле должно получаться промежуточное количество. В этой картине излученная энергия непрерывно зависит от угла.

\begin{figure}[htbp]
\centering
\begin{coursefitfigure}
\begin{tikzpicture}[>=Latex,thick,scale=0.92]
  \draw[rounded corners] (-0.2,-0.2) rectangle (3.0,3.4);
  \node at (1.4,3.05) {\small preparation};
  \draw[->,very thick] (1.4,0.45) -- (1.4,2.55);
  \node[right] at (1.4,2.35) {\small applied field};
  \fill (1.4,1.35) circle (2.3pt);
  \draw[->,very thick] (1.4,1.35) -- (2.25,1.95);
  \node[below] at (1.4,0.15) {\small choose a direction and wait};

  \draw[->] (3.25,1.6) -- (4.25,1.6);

  \draw[rounded corners] (4.5,-0.2) rectangle (8.8,3.4);
  \node at (6.65,3.05) {\small measurement in a fixed direction};
  \draw[->,very thick] (5.15,0.45) -- (5.15,2.55);
  \fill (5.15,1.35) circle (2.3pt);
  \node[anchor=west] at (5.7,2.0) {no photon};
  \node[anchor=west] at (5.7,1.25) {one photon of fixed energy};
\node[anchor=west] at (5.7,0.5) {no intermediate outcome};
\end{tikzpicture}
\end{coursefitfigure}
\caption{Операциональная загадка. За непрерывно варьируемым приготовлением следует детектор с двумя возможными выходами.}
\end{figure}

Для электрона это классическое предсказание неверно. Когда включают поле детектора, возможны только два результата:
\begin{enumerate}
  \item фотон не испускается;
  \item испускается один фотон с фиксированной энергией, соответствующей промежутку между нижним и верхним уровнями.
\end{enumerate}
Никогда не бывает фотона, несущего произвольную долю этой энергии. Непрерывный классический отклик здесь свелся к двоичному квантовому ответу.

Направление, использованное при приготовлении, не стало несущественным. Оно определяет \emph{вероятности} двух ответов. Если приготовить электрон вверх относительно детектора, результат «вверх» получается наверняка. Если приготовить его вниз, переход происходит наверняка. Если приготовить его горизонтально, то вертикальный детектор дает оба ответа одинаково часто:
\begin{equation}
  P(\text{up})=\frac12,
  \qquad
  P(\text{down})=\frac12.
  \label{eq:l02-horizontal-probabilities}
\end{equation}
Таким образом, непрерывно меняющееся семейство приготовлений приводит к двухзначному измерению, статистика которого изменяется непрерывно.

\begin{classroomqa}
\lecturetimestamp{00:21:23}
\audiencequestion{Может ли другой маленький объект, скажем второй электрон, создавать магнитное поле вместо макроскопического магнита?}
\lecturerresponse{Может, но тогда два электрона образуют совместную квантовую систему и могут запутаться. Именно к этому и ведет курс. Мы отложим это, пока не разберемся с пространством состояний одного электрона.}
\end{classroomqa}

\begin{classroomqa}
\lecturetimestamp{00:22:17}
\audiencequestion{Нужно ли находить распределение вероятностей, повторяя весь цикл приготовления и измерения? Что произойдет, если снова измерить тот же электрон?}
\lecturerresponse{Да, исходное распределение требует многократных испытаний с одинаково приготовленными системами. После первого измерения повторение того же измерения для этого электрона дает тот же результат; измерение приготовило его относительно данного детектора. Ожидание излучательного перехода само по себе — одна из возможных процедур приготовления.}
\end{classroomqa}

\begin{classroomqa}
\lecturetimestamp{00:23:48}
\audiencequestion{Представляем ли мы себе свободный электрон, и остается ли магнитное поле все время включенным?}
\lecturerresponse{Мы выделяем спиновую степень свободы и игнорируем положение электрона, независимо от того, свободно ли он движется или как-то иначе ограничен. Большой электромагнит можно быстро включить для приготовления или детектирования. Упрощенный эксперимент устроен так, чтобы задавать вопросы только о спине.}
\end{classroomqa}

\begin{classroomqa}
\lecturetimestamp{00:25:22}
\audiencequestion{Если электрон приготовлен точно в состоянии «вниз» и поле обращают, испускается ли фотон наверняка, и откуда берется его энергия?}
\lecturerresponse{Для точной противоположной ориентации переход происходит с вероятностью единица. При наклонной подготовке вероятность меньше. Энергия в конечном счете берется из работы, совершаемой при создании или обращении внешнего магнитного поля; установка входит в энергетический баланс.}
\end{classroomqa}

\begin{classroomqa}
\lecturetimestamp{00:27:31}
\audiencequestion{Когда появляется фотон? Мгновенен ли переход?}
\lecturerresponse{Время испускания случайно, как и время распада неустойчивого состояния. Можно указать характерное время жизни или период полураспада, и более сильное поле сокращает этот масштаб времени. В более тонкое описание входит соотношение неопределенностей энергия--время, но сейчас существенны лишь два возможных выхода и их измеренные частоты.}
\end{classroomqa}

\begin{classroomqa}
\lecturetimestamp{00:29:17}
\audiencequestion{Не свелся бы пучок, содержащий миллион одинаково приготовленных электронов, просто к измеримой интенсивности света?}
\lecturerresponse{Да. Интенсивность по ансамблю выявляет среднюю вероятность перехода, тогда как время испускания любого отдельного электрона остается случайным. Ансамбль показывает распределение, которое одно испытание показать не может.}
\end{classroomqa}

Это и есть отправная точка. Две изолированные классические точки могут кодировать ответы детектора, но они не могут кодировать все приготовления, которые изменяют вероятности этих ответов. Нам нужно другое пространство состояний.

\section{Комплексные числа и геометрия пространства состояний}

Требуемое пространство состояний линейно: его элементы можно складывать и умножать на скаляры. Эти скаляры — комплексные числа. Запишем
\begin{equation}
  z=a+ib,
  \qquad
  i^2=-1,
  \qquad
  z^*=a-ib,
  \label{eq:l02-complex-number}
\end{equation}
где \(a\) и \(b\) — вещественные числа. Комплексное сопряжение дает положительную величину
\begin{align}
  z^*z
  &=(a-ib)(a+ib) \notag\\
  &=a^2+b^2.
  \label{eq:l02-complex-norm}
\end{align}
Это простое сокращение — алгебраическое зерно правила вероятностей.

\begin{figure}[htbp]
\centering
\includegraphics[width=0.74\linewidth]{lecture_02_figure_02.png}
\caption{Комплексное сопряжение и \(z^*z\) разбираются на доске до того, как амплитуды интерпретируются как вероятности.}
\lectureframe{00:37:10}
\end{figure}

Комплексное векторное пространство содержит объекты \(\ket a,\ket b,\ldots\), такие что при любом комплексном \(c\) выражения
\begin{equation}
  c\ket a,
  \qquad
  \ket a+\ket b
\end{equation}
снова являются векторами. Столбцовые векторы дают конкретную реализацию:
\begin{equation}
  \begin{aligned}
    \ket a&=
    \begin{pmatrix}
      a_1\\a_2\\ \vdots\\a_n
    \end{pmatrix},
    &
    c\ket a&=
    \begin{pmatrix}
      ca_1\\ca_2\\ \vdots\\ca_n
    \end{pmatrix},\\[0.7em]
    \ket a+\ket b&=
    \begin{pmatrix}
      a_1+b_1\\a_2+b_2\\ \vdots\\a_n+b_n
    \end{pmatrix}.
  \end{aligned}
  \label{eq:l02-vector-operations}
\end{equation}
Даже сами комплексные числа образуют одномерное комплексное векторное пространство. Для спина электрона нам понадобится двумерное пространство.

Каждому столбцовому кету сопоставим сопряженный бра-вектор в виде строки:
\begin{equation}
  \ket a=
  \begin{pmatrix}a_1\\a_2\end{pmatrix}
  \quad\longleftrightarrow\quad
  \bra a=
  \begin{pmatrix}a_1^*&a_2^*\end{pmatrix}.
  \label{eq:l02-bra-ket-components}
\end{equation}
Тогда скалярное произведение равно
\begin{equation}
  \braket{b}{a}=b_1^*a_1+b_2^*a_2,
  \label{eq:l02-inner-product}
\end{equation}
а квадрат нормы равен
\begin{equation}
  \braket{a}{a}=|a_1|^2+|a_2|^2.
  \label{eq:l02-norm}
\end{equation}
Это вещественная и неотрицательная величина, потому что каждый член имеет вид \(z^*z\).

\begin{classroomqa}
\lecturetimestamp{00:45:24}
\audiencequestion{Почему у символа бра нет явной звездочки?}
\lecturerresponse{Преобразование кета в его бра уже означает транспонирование и комплексное сопряжение. Разворот скобки сам по себе обозначает эту операцию, так что запись еще одной звездочки просто повторяла бы то, что уже подразумевается.}
\end{classroomqa}

\section{Базис кубита и правило Борна}

Вернемся теперь к детектору. Пусть \(\ket{+}\) обозначает состояние «вверх», а \(\ket{-}\) — состояние «вниз». Общее спиновое состояние является линейной комбинацией
\begin{equation}
  \ket\psi=a_+\ket{+}+a_-\ket{-}
  \quad\longleftrightarrow\quad
  \begin{pmatrix}a_+\\a_-\end{pmatrix}.
  \label{eq:l02-general-qubit}
\end{equation}
Знаки плюс и минус внутри кетов — это метки, а не арифметические операции. После выбора этого упорядоченного базиса
\begin{equation}
  \ket{+}=
  \begin{pmatrix}1\\0\end{pmatrix},
  \qquad
  \ket{-}=
  \begin{pmatrix}0\\1\end{pmatrix}.
  \label{eq:l02-basis-columns}
\end{equation}
В самом деле,
\begin{equation}
  a_+\begin{pmatrix}1\\0\end{pmatrix}
  +a_-\begin{pmatrix}0\\1\end{pmatrix}
  =
  \begin{pmatrix}a_+\\a_-\end{pmatrix}.
  \label{eq:l02-basis-expansion}
\end{equation}

\begin{figure}[htbp]
\centering
\includegraphics[width=0.72\linewidth]{lecture_02_figure_03.png}
\caption{Состояния детектора задаются двумя стандартными базисными столбцами.}
\lectureframe{00:58:10}
\end{figure}

Коэффициенты — это амплитуды вероятности. Вероятность каждого исхода детектора равна квадрату модуля соответствующей амплитуды:
\begin{equation}
  P(+)=|a_+|^2,
  \qquad
  P(-)=|a_-|^2.
  \label{eq:l02-born-rule}
\end{equation}
Поскольку это единственные исходы,
\begin{equation}
  \begin{gathered}
    |a_+|^2+|a_-|^2=1,\\
    \text{or equivalently}\qquad \braket{\psi}{\psi}=1.
  \end{gathered}
  \label{eq:l02-normalization}
\end{equation}
Следовательно, физический вектор состояния нормирован.

\begin{classroomqa}
\lecturetimestamp{00:56:54}
\audiencequestion{Имеют ли знаки плюс и минус алгебраический смысл в компонентах?}
\lecturerresponse{Здесь они обозначают два базисных состояния и их коэффициенты. Мы могли бы назвать их \(1\) и \(2\) или «вверх» и «вниз». В частности, \(-\ket{+}\) является допустимым скалярным кратным \(\ket{+}\); это не состояние \(\ket{-}\).}
\end{classroomqa}

Умножение состояния на комплексное число единичного модуля не меняет ни одной из его вероятностей:
\begin{equation}
  \ket\psi\longmapsto e^{i\phi}\ket\psi,
  \qquad
  |e^{i\phi}a_\pm|^2=|a_\pm|^2.
  \label{eq:l02-global-phase}
\end{equation}
Следовательно, два вектора, различающиеся только общей фазой, представляют одно и то же физическое состояние. \emph{Относительная} фаза — другое дело. Например, если заменить только \(a_-\) на \(-a_-\), то это, вообще говоря, нельзя устранить умножением всего вектора на одну фазу; это может изменить результаты измерений в другом базисе.

\begin{classroomqa}
\lecturetimestamp{00:59:54}
\audiencequestion{Являются ли \(\ket{+}\), \(-\ket{+}\) и \(e^{i\phi}\ket{+}\) разными состояниями? А если изменить знак одного члена в суперпозиции?}
\lecturerresponse{Это разные математические векторы, но, когда этот множитель относится ко всему вектору и имеет единичный модуль, они описывают одно и то же физическое состояние. Относительный знак или фаза между членами «вверх» и «вниз» наблюдаемы при подходящих измерениях и поэтому меняют состояние.}
\end{classroomqa}

Базис ортонормирован:
\begin{align}
  \braket{+}{+}&=1,
  &\braket{-}{-}&=1,
  \notag\\
  \braket{+}{-}&=0,
  &\braket{-}{+}&=0.
  \label{eq:l02-orthonormality}
\end{align}
Например,
\begin{equation}
  \braket{+}{-}
  =
  \begin{pmatrix}1&0\end{pmatrix}
  \begin{pmatrix}0\\1\end{pmatrix}
  =0.
\end{equation}

\begin{figure}[htbp]
\centering
\includegraphics[width=0.76\linewidth]{lecture_02_figure_04.png}
\caption{Четыре скалярных произведения базиса «вверх/вниз» непосредственно проверяются на доске.}
\lectureframe{01:05:00}
\end{figure}

Состояние с равными весами даёт первую полезную поправку на доске. Столбец
\begin{equation}
  \begin{pmatrix}1\\1\end{pmatrix}
\end{equation}
имеет квадрат нормы \(2\), а не \(1\), и потому не нормирован. Деление на \(\sqrt2\) исправляет это:
\begin{equation}
  \ket{\rightarrow}
  =
  \frac{1}{\sqrt2}
  \begin{pmatrix}1\\1\end{pmatrix}
  =
  \frac{\ket{+}+\ket{-}}{\sqrt2}.
  \label{eq:l02-horizontal-state}
\end{equation}
Оно даёт \(P(+)=P(-)=1/2\) в вертикальном детекторе и соответствует горизонтальной подготовке. Исправление от \((1,1)^T\) к \(2^{-1/2}(1,1)^T\) не косметическое: именно нормировка превращает квадраты амплитуд в полное распределение вероятностей.

После перерыва стоит вернуться к двум значениям слова «вектор». Пространственная стрелка имеет три вещественные компоненты. Вектор состояния кубита имеет две комплексные компоненты, то есть четыре вещественных числа до наложения нормировки и отождествления общей фазы. Эти объекты связаны, но один не сводится просто к переобозначению другого. Полное геометрическое соответствие между ними нам пока не нужно.

\begin{classroomqa}
\lecturetimestamp{01:10:56}
\audiencequestion{Откуда мы знаем размерность пространства состояний для другой квантовой системы?}
\lecturerresponse{Эксперимент показывает, сколько взаимно исключающих исходов может различить полное измерение. Спин электрона имеет два таких исхода и потому ему соответствует двумерное комплексное пространство состояний. Система со спином один имеет три исхода в соответствующем измерении и требует трёх компонент.}
\end{classroomqa}

\section{Классические наблюдаемые и ансамблевые средние}

Описав состояния, перейдём к измеримым величинам. В классической системе с дискретными состояниями \(n\) наблюдаемая — это принимающая вещественные значения функция \(f_n\), заданная на этих состояниях. Если состояние неопределённо и задаётся вероятностями \(p_n\), её среднее по ансамблю равно
\begin{equation}
  \langle f\rangle_{\mathrm{ens}}
  =\sum_n p_n f_n.
  \label{eq:l02-classical-expectation}
\end{equation}
Эта запись не означает, что какое-либо отдельное измерение вообще должно быть равно этому числу.

Для монеты зададим \(f_+=+1\) и \(f_-=-1\). Тогда
\begin{equation}
  \langle f\rangle_{\mathrm{ens}}
  =p_+-p_-.
  \label{eq:l02-coin-expectation}
\end{equation}
У честной монеты среднее равно нулю. Если \(p_+=3/4\) и \(p_-=1/4\), то среднее равно \(1/2\). Ни \(0\), ни \(1/2\) не являются исходом отдельного броска; каждое из них обобщает множество результатов \(+1\) и \(-1\).

\begin{classroomqa}
\lecturetimestamp{01:23:30}
\audiencequestion{Как ожидаемое значение может быть равно нулю или одной второй, если ни одно из этих значений не может появиться в отдельном испытании?}
\lecturerresponse{Выражение «ожидаемое значение» потенциально вводит в заблуждение. Оно означает среднее по большому ансамблю, а не значение, которое следует ожидать в следующем запуске. Для честной последовательности результатов \(+1\) и \(-1\) среднее равно нулю, а при соотношении три к одному — одной второй.}
\end{classroomqa}

Особенно полезная наблюдаемая — индикатор одного состояния:
\begin{equation}
  \chi_k(n)=
  \begin{cases}
    1,&n=k,\\
    0,&n\ne k.
  \end{cases}
\end{equation}
Его среднее равно
\begin{equation}
  \langle\chi_k\rangle_{\mathrm{ens}}=p_k.
\end{equation}
Следовательно, знание средних значений всех классических наблюдаемых включает, в частности, знание распределения вероятностей.

\section{Матрицы и квантовые наблюдаемые}

В квантовой механике наблюдаемые представляются не функциями на точках пространства состояний, а линейными операторами. В конечномерном пространстве состояний эти операторы являются матрицами. Матрица \(M\) отображает один вектор в другой:
\begin{equation}
  M\ket a=\ket{a'}.
  \label{eq:l02-matrix-action}
\end{equation}
В случае двух компонент,
\begin{equation}
  \begin{pmatrix}
    m_{11}&m_{12}\\
    m_{21}&m_{22}
  \end{pmatrix}
  \begin{pmatrix}a_1\\a_2\end{pmatrix}
  =
  \begin{pmatrix}
    m_{11}a_1+m_{12}a_2\\
    m_{21}a_1+m_{22}a_2
  \end{pmatrix}.
  \label{eq:l02-matrix-components}
\end{equation}

\begin{figure}[htbp]
\centering
\includegraphics[width=0.78\linewidth]{lecture_02_figure_05.png}
\caption{Действие матрицы построчно расписывается в две компоненты преобразованного вектора.}
\lectureframe{01:32:10}
\end{figure}

На вещественных двумерных векторах геометрия становится наглядной. Матрицы
\begin{equation}
  \begin{pmatrix}2&0\\0&2\end{pmatrix},
  \qquad
  \begin{pmatrix}1&0\\0&2\end{pmatrix},
  \qquad
  \begin{pmatrix}2&0\\0&1\end{pmatrix}
\end{equation}
соответственно одинаково растягивают все направления, растягивают только \(y\)-компоненту и растягивают только \(x\)-компоненту. Матрица
\begin{equation}
  \begin{pmatrix}0&1\\1&0\end{pmatrix}
\end{equation}
меняет компоненты местами и отражает плоскость относительно \(x=y\). А
\begin{equation}
  \begin{pmatrix}0&1\\-1&0\end{pmatrix}
  \begin{pmatrix}v_x\\v_y\end{pmatrix}
  =
  \begin{pmatrix}v_y\\-v_x\end{pmatrix}
  \label{eq:l02-quarter-turn}
\end{equation}
поворачивает векторы на четверть оборота по часовой стрелке в той конвенции, которая использовалась на доске. Косой сдвиг даёт ещё один линейный пример, хотя его матрица оставлена в качестве упражнения. Существенное ограничение состоит в линейности: не всякое мыслимое преобразование плоскости представимо матрицей.

Квантовые наблюдаемые принадлежат к особому классу.

\begin{definition}
Матрица \(M\) называется \emph{эрмитовой}, если
\begin{equation}
  M_{ij}=M_{ji}^*.
  \label{eq:l02-hermitian-condition}
\end{equation}
\end{definition}

Диагональное условие \(M_{ii}=M_{ii}^*\) требует, чтобы каждый диагональный элемент был вещественным. Внедиагональные элементы образуют пары комплексно-сопряжённых чисел. Например,
\begin{equation}
  M=
  \begin{pmatrix}
    r&4+i\\
    4-i&s
  \end{pmatrix},
  \qquad r,s\in\mathbb R,
  \label{eq:l02-hermitian-example}
\end{equation}
является эрмитовой матрицей. Если все элементы оказываются вещественными, то эрмитовость совпадает с симметричностью.

\begin{figure}[htbp]
\centering
\includegraphics[width=0.75\linewidth]{lecture_02_figure_06.png}
\caption{Определяющее сопряжение относительно диагонали, \(M_{ij}=M_{ji}^*\), отдельно выделено в конце лекции.}
\lectureframe{01:42:50}
\end{figure}

Эрмитовы матрицы будут представлять измеримые величины, включая спин вдоль направлений \(x\), \(y\) и \(z\). Их не следует путать с операторами, описывающими физический толчок или временную эволюцию.

\begin{classroomqa}
\lecturetimestamp{01:45:49}
\audiencequestion{Если матрица действует на состояние как толчок, переводя его в другое состояние, должна ли она сохранять нормировку? Сохраняет ли её эрмитова матрица?}
\lecturerresponse{Физический толчок должен сохранять длину вектора состояния, но произвольная эрмитова матрица этого не делает. Преобразования, сохраняющие длину, унитарны. Эрмитовы операторы представляют наблюдаемые; унитарные операторы представляют обратимые преобразования, действующие на состояния. Их подробная связь будет рассмотрена далее.}
\end{classroomqa}

Мы дошли до минимального языка, необходимого для следующего шага. Кубит — это нормированный вектор в двумерном комплексном векторном пространстве; амплитуды определяют вероятности; а наблюдаемые представляются эрмитовыми матрицами. Имея эти элементы, можно обсуждать спин вдоль произвольных направлений и, в конечном счёте, запутанные пары, не возвращаясь к классической картине, которая уже оказалась несостоятельной.
